% Lösungen Einheit 5 von 5 – Ortskurve und Kurvenvergleich (zusammenbau v0.1)
\einheitenkopf{Einheit 5 von 5 · Ortskurve und Kurvenvergleich}

\erg{18}{a) trägt den Parameter \quad b) Scheitel $(a \,|\, -a^2)$; aus $x = a$ folgt $y = -x^2$. \quad c) $f_2(x) = x^2 - 4x + 3$, Tiefpunkt $(2 | -1)$; $3 - 2^2 = -1$: Der Punkt liegt auf der Kurve. \quad d) $m = \frac{9 - 5}{4 - 2} = 2$; Ortsgerade $y = 2x + 1$ (nicht der Quotient der Koordinaten eines Punkts). \quad e) $f_a\left(\frac{3}{a}\right) = \frac{27}{a^3} \cdot e^{-3}$ und $\frac{1}{e^3} \cdot \left(\frac{3}{a}\right)^3 = \frac{27}{a^3 e^3}$: gleich für jedes $a$. \quad f) Hochpunkte $(k \,|\, 2k)$ und $(k + 1 \,|\, 2k + 2)$; Abstand $\sqrt{1^2 + 2^2} = \sqrt{5} \approx 2{,}24$ für jedes $k$. \quad g) $-f(x) = \frac{1}{4}x(x - 5)e^{x} = g_5(x)$, also $b = 5$. \quad h) $k x^2 = 2k x$ liefert $x = 0$ oder $x = 2$: gemeinsame Punkte $(0 | 0)$ und $(2 \,|\, 4k)$. \quad i) Sonderfall: $f_1'(x) = \frac{3}{2}(1 - x^2)$, Extrempunkte $(1 | 1)$ und $(-1 | -1)$, beide auf $y = x$. Die Streckung mit gleichem Faktor in beiden Richtungen führt Extrempunkte in Extrempunkte über und bildet die Ursprungsgerade $y = x$ auf sich ab: Alle Extrempunkte liegen auf $y = x$.}
\erg{19}{a) Paul hat $x$ für $a$ eingesetzt, statt $a$ durch $x$ auszudrücken. Richtig: $a = \frac{x}{2}$, also $y = -4 \cdot \frac{x^2}{4} = -x^2$. \quad b) Für jedes $a$ erfüllen $x = a$ und $y = 2a^2$ die Gleichung $y = 2x^2$; da $a$ beliebig war, liegt jeder Extrempunkt darauf.}
